%%%% ijcai23.tex

\typeout{IJCAI--23 Instructions for Authors}

% These are the instructions for authors for IJCAI-23.

\documentclass{article}
\pdfpagewidth=8.5in
\pdfpageheight=11in

% The file ijcai23.sty is a copy from ijcai22.sty
% The file ijcai22.sty is NOT the same as previous years'
\usepackage{ijcai23}

% Use the postscript times font!
\usepackage{times}
\usepackage{soul}
\usepackage{url}
\usepackage[hidelinks]{hyperref}
\usepackage[utf8]{inputenc}
\usepackage[small]{caption}
\usepackage{graphicx}
\usepackage{amsmath}
\usepackage{amsthm}
\usepackage{booktabs}
\usepackage{algorithm}
\usepackage{algorithmic}
\usepackage[switch]{lineno}
\setcounter{secnumdepth}{3}
\usepackage{multirow}
\usepackage[table,xcdraw]{xcolor}
\definecolor{myurlcolor}{rgb}{0,0.5,0.5}  % Custom color.
\usepackage{tabu}
\newcommand{\xz}[1]{\textcolor{black}{#1}}

\hypersetup{
  colorlinks=true,  % Colors the text of the links and anchors.
  urlcolor=myurlcolor,  % Apply custom color only to URLs.
  citecolor=black,  % Keeps citation links black (or any color you prefer).
  linkcolor=black,  % Keeps internal links like references black (or any color you prefer).
}

\newenvironment{subgroup}
  {$\left\{\tabular{l}}
  {\endtabular\right.$}
% \usepackage[colorlinks,linkcolor=blue]{hyperref}
% \usepackage{cite}
% Comment out this line in the camera-ready submission
% \linenumbers

\urlstyle{same}

% the following package is optional:
%\usepackage{latexsym}

% See https://www.overleaf.com/learn/latex/theorems_and_proofs
% for a nice explanation of how to define new theorems, but keep
% in mind that the amsthm package is already included in this
% template and that you must *not* alter the styling.
\newtheorem{example}{Example}
\newtheorem{theorem}{Theorem}

% Following comment is from ijcai97-submit.tex:
% The preparation of these files was supported by Schlumberger Palo Alto
% Research, AT\&T Bell Laboratories, and Morgan Kaufmann Publishers.
% Shirley Jowell, of Morgan Kaufmann Publishers, and Peter F.
% Patel-Schneider, of AT\&T Bell Laboratories collaborated on their
% preparation.

% These instructions can be modified and used in other conferences as long
% as credit to the authors and supporting agencies is retained, this notice
% is not changed, and further modification or reuse is not restricted.
% Neither Shirley Jowell nor Peter F. Patel-Schneider can be listed as
% contacts for providing assistance without their prior permission.

% To use for other conferences, change references to files and the
% conference appropriate and use other authors, contacts, publishers, and
% organizations.
% Also change the deadline and address for returning papers and the length and
% page charge instructions.
% Put where the files are available in the appropriate places.


% PDF Info Is REQUIRED.
% Please **do not** include Title and Author information
\pdfinfo{
/TemplateVersion (IJCAI.2023.0)
}

\title{Large Language Model based Multi-Agents: A Survey of Progress and Challenges}


% Single author syntax
% \author{
%     Author Name
%     \affiliations
%     Affiliation
%     \emails
%     email@example.com
% }

% Multiple author syntax (remove the single-author syntax above and the \iffalse ... \fi here)
% \iffalse
\newcommand*\samethanks[1][\value{footnote}]{\footnotemark[#1]}
\author{
Taicheng Guo$^1$
\and
Xiuying Chen$^2$\and
Yaqi Wang$^3$\thanks{This work was done when Yaqi was visiting students at the University of Notre Dame.} \and
Ruidi Chang \and 
Shichao Pei$^4$ \and \\
Nitesh V. Chawla$^1$ \and
Olaf Wiest$^1$ \and
Xiangliang Zhang$^1$\thanks{Corresponding author.}
\affiliations
$^1$University of Notre Dame, $^2$King Abdullah University of Science and Technology\\
$^3$Southern University of Science and Technology,
$^4$University of Massachusetts Boston
\emails
\{tguo2, nchawla, owiest, xzhang33\}@nd.edu,
xiuying.chen@kaust.edu.sa,
ywang84@nd.edu, ruidic@alumni.cmu.edu, shichao.pei@umb.edu
}
% \fi

\begin{document}

\maketitle

\begin{abstract}
% 1. Introduction of LLM and Agents
Large Language Models (LLMs) have achieved remarkable success across a wide array of tasks. Due to the impressive planning and reasoning abilities of LLMs, they have been used as autonomous agents to do many tasks automatically. Recently, based on the development of using one LLM as a single planning or decision-making agent, LLM-based multi-agent systems have achieved considerable progress in complex problem-solving and world simulation. 
To provide the community with an overview of this dynamic field,  we present this survey to offer an in-depth discussion on the essential aspects of multi-agent systems based on LLMs, as well as the challenges. Our goal is for readers to gain substantial insights on the following questions: What domains and environments do LLM-based multi-agents simulate? How are these agents profiled and how do they communicate? What mechanisms contribute to the growth of agents' capacities? For those interested in delving into this field of study, we also summarize the commonly used datasets or benchmarks for them to have convenient access. To keep researchers updated on the latest studies,  we maintain an open-source \href{https://github.com/taichengguo/LLM_MultiAgents_Survey_Papers}{GitHub repository}, dedicated to outlining the research on LLM-based multi-agent systems. %(\href{https://github.com/taichengguo/LLM_MultiAgents_Survey_Papers}{LLM\_MultiAgents\_Survey\_Papers}).  

% 2. The structure of our paper 
%In this survey, we systematically review the progress of LLM-based multi-agent studies. Specifically, we first introduce the basic concepts of the LLM-based multi-agent system. Then we organize and categorize existing works into two main categories based on their motivation and the two main categories contain several sub-categories based on the practical scenario. We also highlight the challenges and opportunities of this promising technique. 
% 3. Open resources
%Finally, We release an open-source repository \href{https://github.com/taichengguo/LLM_MultiAgents_Survey_Papers}{LLM\_MultiAgents\_Survey\_Papers} which includes constantly updated LLM based Multi-Agents papers, providing a useful resource for various communities.
\end{abstract}

\section{Introduction}

% 1. Impact of Large lanugage model: 
    % Revolutionize human-like intelligence
% 2. LLM-based agent: 
    % Why LLM-based agent; to maximize the intelligence, LLM-based agent is proposed and developed.
    % What can LLM-based agent do / The importance of LLM-based Agent

% Impact of Large lanugage model: 
Large Language Models (LLMs) have recently shown remarkable potential in reaching a level of reasoning and planning capabilities comparable to humans.
% Why LLM-based agent;
This ability exactly aligns with the expectations of humans for autonomous agents that can perceive the surroundings, make decisions, and take actions in response~\cite{xi2023rise,Wooldridge1995IntelligentAT,Artificial_Intelligence_AModernApproach,guo2023indeed,liang2023let}. Hence, LLM-based agent has been studied and rapidly developed to understand and generate human-like instructions, facilitating sophisticated interactions and decision-making in a wide range of contexts
%to maximize this potential of LLMs, LLM-based agent has been studied and rapidly developed. 
% What is LLM-based agent
%The LLM-based agent refers to an autonomous agent powered by the LLM, enabling it to understand and generate human-like instructions, facilitating sophisticated interactions and decision-making in a wide range of contexts.
% LLM-based agent: previous work
%The LLM-based agent topic has attracted a lot of interests and various works have been developed by researchers
~\cite{yao2023tree,shinn2023reflexion,li2023apibank}. Timely survey papers systematically summarize the progress of LLM-based agents, as seen in works ~\cite{xi2023rise,wang2023survey}.
%Several systematic survey works~\cite{xi2023rise,wang2023survey} were also conducted to summarize the LLM-based agent. 


% \begin{figure}[tp]
% \begin{flushleft} % This will align the subgroup and caption to the left
% LLM-based Multi-Agents
% \begin{subgroup}
%   Idea Principal \\[2em]
%   Idea Principal \\[2em]
%   Idea Principal
% \end{subgroup}
% \caption{Caption describing the relationship between the general idea and the main ideas.}
% \end{flushleft}
% \end{figure}

\begin{figure*}[ht]
\begin{center}
\includegraphics[scale=0.56]{img/Trend_example.png}
\end{center}
\caption{The rising trend in the research field of LLM-based Multi-Agents. For Problem Solving and World Simulation, we categorize current work into several categories and count the number of papers of different types at 3-month intervals. The number at each leaf node denotes the count of papers within that category.}
\label{trend}
\end{figure*}


% 3. LLM-based Multi-Agent:
    % Why we need Multi-Agent
    % What is Multi-Agent
    % Why we focus on the Multi-Agent (importance of multi-agent)

% Why we need LLM-based Multi-Agent: Collective Intelligence
Based on the inspiring capabilities of the single LLM-based agent, LLM-based Multi-Agents have been proposed to leverage the collective intelligence and specialized profiles and skills of multiple agents. 
% What is LLM-based Multi-Agent 
\xz{Compared to systems using a single LLM-powered agent, multi-agent systems offer advanced capabilities by \underline{1)} specializing LLMs into various distinct agents, each with different capabilities, and \underline{2)} enabling interactions among these diverse agents to simulate complex real-world environments effectively. In this context, multiple autonomous agents collaboratively engage in planning, discussions, and decision-making, mirroring the cooperative nature of human group work in problem-solving tasks. This approach capitalizes on the communicative capabilities of LLMs, leveraging their ability to generate text for communication and respond to textual inputs. Furthermore, it exploits LLMs' extensive knowledge across various domains and their latent potential to specialize in specific tasks. Recent research has demonstrated promising results in utilizing LLM-based multi-agents for solving various tasks,}
%The LLM-based Multi-Agents refers to a system where multiple autonomous agents are empowered by LLM. These agents utilize the capabilities of LLM to interact with each other or the environment of the system. 
% Why we focus on the Multi-Agent (importance of multi-agent, research trend: not AI but interdisciplinary researchers => This Multi-Agent have broader impact to many fields)
%Compared to the individual agent, in the LLM-based Multi-Agents system, each agent's specialized capabilities contribute to the system's ability to efficiently scale up and adapt to changing environments 
such as software development~\cite{hong2023metagpt,qian2023communicative}, multi-robot systems~\cite{mandi2023roco,zhang2023building}, %etc. Additionally, the communication between agents in the Multi-Agents system enables the simulation of complex scenarios such as 
society simulation~\cite{park2023generative,park2022social}, policy simulation~\cite{xiao2023simulating,hua2023war}, and game simulation~\cite{xu2023language,wang2023avalon}. 
\xz{Due to the nature of interdisciplinary study in this field, it has attracted a diverse range of researchers, expanding beyond AI experts to include those from social science, psychology, and policy research. The volume of research papers is rapidly increasing,}
%Attracted by the powerful capabilities of the LLM-based Multi-agents systems, an increasing variety of researchers, extending beyond just AI researchers to those in fields such as social science, psychology, and policy researchers have delved into this promising research direction 
as shown in Fig. \ref{trend} (inspired by the design in ~\cite{gao2023retrieval}),
%\footnote{\scriptsize The way our figure is designed is also inspired by ~\cite{gao2023retrieval}.}, 
thus broadening the impact of LLM-based Multi-Agent research.
% Limitations (No summary)
Nonetheless, earlier efforts were undertaken independently,  resulting in an absence of a systematic review to summarize them, establish comprehensive \xz{blueprint of this field}, and examine future research challenges. This underscores the significance of our work and serves as the motivation behind presenting this survey paper, dedicated to the research on LLM-based multi-agent systems.


% 4. Objective of this survey:
    % Our purpose and contribution: AI researchers, Interdisciplinary researchers
      % how our survey aligns with current trends in AI and LLM research.

We expect that our survey can make significant contributions to both the research and development of LLMs and to a wider range of interdisciplinary studies employing LLMs. Readers will gain a comprehensive overview of  LLM-based Multi-Agent (LLM-MA) systems, grasp the fundamental concepts involved in establishing multi-agent systems based on LLMs, and catch the latest research trends and applications in this dynamic field. We recognize that this field is in its early stages and is rapidly evolving with fresh methodologies and applications. To provide a sustainable resource complementing our survey paper, we maintain an open-source GitHub repository\footnote{\scriptsize \url{https://github.com/taichengguo/LLM_MultiAgents_Survey_Papers}}.
We hope that our survey will inspire further exploration and innovation in this field,
%LLM-based Multi-Agent techniques and 
as well as applications across a wide array of research disciplines.

%Given the tremendous effort in this rapidly developing research area, we are motivated to present a comprehensive and systematic survey focusing on the LLM-based Multi-Agents. 
% Our purpose and contribution: AI researchers, Interdisciplinary researchers
%We expect that our survey can contribute to both LLM research fields and the broader interdisciplinary studies of LLM-based Multi-Agents, encompassing areas such as society, psychology, and policy, etc. For researchers who are already familiar with LLM-based Multi-Agents concepts,  our survey presents the latest applications to demonstrate current research trends. For newcomers or Non-CS researchers who are interested in the use of LLM-based Multi-Agents for problem solving or world simulation, we provide systematic taxonomy of LLM-based Multi-Agents to help them obtain a quick grasp of the core concepts of LLM-based Multi-Agents.
% Potential future Challenges and developments
%We are the view that this field is still in its early phase and we hope that our survey will motivate further exploration and innovation in LLM-based Multi-Agent techniques and applications across a diverse spectrum of research disciplines.

% 5. The detail structure of the survey paper
% The specific aspects of our survey: 
% How we organize the current work: Scope / Taxonomy, ......     Potential future Challenges and developments
%We organize this survey in the following structure. Related works are presented in Section~\ref{Sec2:background}. By drawing connections among the existing LLM-based Multi-Agents research, we first provide a detailed taxonomy, which summarizes four key properties (Agents-Environment Interface, Agents Profiling, Agents Communication and Agents Capabilities Acquisition) of the LLM-based Multi-Agents systems in Section~\ref{Sec3:Taxonomy}. Then we summarize the LLM-based Multi-Agents Framework in Section~\ref{Sec4:Framework} and Multi-Agents Orchestration and Efficiency in Section~\ref{Sec5:Orchestration}. We categorize current applications into two main streams: Multi-Agents for problem solving and Multi-Agents for World Simulation in Section~\ref{Sec5:Application}. We then look into the current LLM-based Multi-Agents datasets and benchmarks in Section~\ref{Sec6:Datasets and Benchmarks}. Based on the previous summary, we open the discussion for future research challenges and opportunities in Section~\ref{Sec7:Challenges and Opportunities}. The conclusions are summarized in Section~\ref{Sec8:Conclusion}.

To assist individuals from various backgrounds in understanding LLM-MA techniques and to complement existing surveys by tackling unresolved questions, we have organized our survey paper in the following manner. After laying out the background knowledge in Section~\ref{Sec2:background}, we address a pivotal question: \emph{How are LLM-MA systems aligned with the collaborative task-solving environment}?   To answer this, we present a comprehensive schema for positioning, differentiating, and connecting various aspects of LLM-MA systems in  \textbf{Section~\ref{Sec3:Taxonomy}}. We delve into this question by discussing: 1) the  \textbf{agents-environment interface}, which details how agents interact with the task environment; 2)  \textbf{agent profiling}, which explains how an agent is characterized by an LLM to behave in specific ways; 3)  \textbf{agent communication}, which examines how agents exchange messages and collaborate; and 4)  \textbf{agent capability acquisition}, which explores how agents develop their abilities to effectively solve problems. 
An additional perspective for reviewing studies about LLM-MA is their application. In  \textbf{Section~\ref{Sec5:Application}}, we categorize current  \textbf{applications} into two primary streams: multi-agents for \textbf{problem-solving} and multi-agents for \textbf{world simulation}.
To guide individuals in identifying appropriate \textbf{tools and resources},  we present open-source implementation frameworks for studying LLM-MA, as well as the usable datasets and benchmarks in  \textbf{Section~\ref{Sec4:Framework}}.
Based on the previous summary, we open the discussion for future research challenges and opportunities in Section~\ref{Sec7:Challenges and Opportunities}. The conclusions are summarized in Section~\ref{Sec8:Conclusion}.
\section{Background}
\label{Sec2:background}
\subsection{Single-Agent Systems Powered LLMs}
%\subsection{LLM-based Single-Agent\protect\footnote{\scriptsize The framework of LLM-based Single-Agent capabilities is also inspired by: \href{https://lilianweng.github.io/posts/2023-06-23-agent}{https://lilianweng.github.io/posts/2023-06-23-agent}}}
% 1. Why we introduce this: Single Agent => Multi-Agents
%Since LLM-based Multi-Agents are composed by a number of LLM-based single agent, we will first introduce key capabilities of LLM-based agent to to provide a foundational understanding of their individual functionalities.
We introduce the background by first outlining the capabilities of a single-agent system based on  LLMs, following the discussion presented in \cite{agent_link}. 

% 2. What is the ability; Representative Work; How it empower single Agent
\paragraph{Decision-making Thought:} 
This term denotes the capability of LLM-based agents, guided by prompts, to break down complex tasks into smaller subgoals~\cite{khot2023decomposed}, think through each part methodically (sometimes exploring multiple paths)~\cite{yao2023tree}, and learn from past experiences~\cite{shinn2023reflexion} to perform better decision-making on complex tasks. This capability enhances the autonomy of a single LLM-based agent and bolsters its effectiveness in problem-solving.

\paragraph{Tool-use:} LLM-based agents' tool-use capability allows them to leverage external tools and resources to accomplish tasks, enhancing their functional capabilities and operate more effectively in diverse and dynamic environments~\cite{li2023apibank,ruan2023tptu,gao2023retrieval}.  

\paragraph{Memory:} This ability refers to the capability of LLM-based agent for conducting in-context learning~\cite{dong2023survey} as short memory or external vector database~\cite{lewis2021retrievalaugmented} as long memory to preserve and retrieve information over prolonged periods~\cite{wang2023survey}. This ability enables a single LLM-based agent to maintain contextual coherence and enhance learning from interactions. 

\subsection{Single-Agent VS. Multi-Agent Systems}
% Strong Single agent capabilities contribute to the multi-agents
Single-Agent systems empowered by LLMs have shown inspiring cognitive abilities \cite{sumers2023cognitive}. 
The construction of such systems concentrates on formulating their internal mechanisms and interactions with the external environment. Conversely, 
%These individual capabilities enable effective communication between agents and thus are essential for the emergence of collective intelligence of LLM-based Multi-Agents.
% Difference between Single Agent and Multi-Agents
%However, unlike the LLM-based single agents focus on internal mechanisms and interactions with the external environment, 
LLM-MA  systems emphasize diverse agent profiles, inter-agent interactions, and collective decision-making processes. From this perspective, more dynamic and complex tasks can be tackled by the collaboration of multiple autonomous agents, each of which is equipped with unique strategies and behaviors, and engaged in communication with one another.



% \subsection{Reinforcement-Learning based Multi-Agents}
% Different from RL Multi-Agents
% The concept of multi-agents is widely discussed in the reinforcement-learning research field. 
\section{Dissecting LLM-MA Systems: Interface, Profiling, Communication, and Capabilities}
\label{Sec3:Taxonomy}
% ddl 11.17
In this section, we delve into the intricacies of LLM-MA systems, where multiple autonomous agents engage in collaborative activities akin to human group dynamics in problem-solving scenarios. \xz{A critical inquiry we address is how these LLM-MA systems are  aligned  to  their operational environments and the collective objectives they are designed to achieve. To shed light on this, we present the general architecture of these systems  in Fig. \ref{fig:Architecture}.} Our analysis dissects the operational framework of these systems, focusing on four key aspects: the agents-environment interface, agent profiling, agent communication, and agent capability acquisition.

\begin{figure*}[ht]
\centering
\includegraphics[width= 0.73\textwidth]{img/LLM-MA.png}
\caption{The Architecture of LLM-MA Systems.} 
\label{fig:Architecture} 
\end{figure*}


%We present a comprehensive taxonomy that offers a wide-ranging perspective of the domain. Concretely, this section introduces the taxonomy of LLM-based Multi-Agents from the following four perspectives: Agents-Environment Interface, Agents Profiling, Agents Communication, and Agents Capabilities. Subsequently, we elaborate on the criteria of the taxonomy in detail.  

\subsection{Agents-Environment Interface}
% 1. What is Environment 
\xz{The operational environments defines the specific contexts or settings in which the LLM-MA systems are deployed and interact. 
For example, these environments can be like software development \cite{hong2023metagpt}, gaming \cite{mao2023alympics}, and various other domains such as financial markets \cite{li2023tradinggpt}  or even social behavior modeling \cite{park2023generative}.  The LLM-based agents perceive and act within the environment, which in turn influences their behavior and decision making. For example, in the Werewolf Game simulation, the sandbox environment sets the game's framework, including transitions from day to night, discussion periods, voting mechanics, and reward rules. Agents, such as werewolves and the Seer, perform specific actions like killing or checking roles. Following these actions, agents receive feedback from the environment, informing them of the game's current state. This information guides the agents in adjusting their strategies over time, responding to the evolving gameplay and interactions with other agents.}
%We begin with the Agents-Environment Interface of the LLM-based Multi-Agents. The Environment of the Multi-agents systems typically refers to the external setting within multiple agents. Agents can interact with the environment to obtain feedback. 
The Agents-Environment Interface refers to the  way in which agents interact with and perceive the environment. It's through this interface that agents understand their surroundings, make decisions, and learn from the outcomes of their actions.
We categorize the current interfaces in LLM-MA systems into three types, \emph{Sandbox, Physcial}, and \emph{None}, as detailed in Table \ref{summary}.
%The Agents-Environment Interface can be categorized into three types: Sandbox, Physcial and None. 
The \emph{\textbf{Sandbox}} refers to a simulated or virtual environment built by human where agents can interact more freely and experiment with various actions and strategies. This kind of interface is widely used in software development (code interpreter as simulated environment)~\cite{hong2023metagpt}, gaming (using game rules as simulated environment)~\cite{mao2023alympics}, etc.  The \emph{\textbf{Physical}} is a real-world environment where agents interact with physical entities and obey real-world physics and constraints. In physical space, agents normally need to take actions that can have direct physical outcomes. For example, in tasks such as sweeping the floor, making sandwiches, packing groceries, and arranging cabinets, robotic agents are required to perform actions iteratively, observe the physical environment, and continuously refine their actions~\cite{mandi2023roco}. Lastly, \emph{\textbf{None}} refers to scenarios where there is no specific external environment, and agents do not interact with any environment. For example, many applications~\cite{du2023improving,xiong2023examining,chan2023chateval} utilize multiple agents to debate a question to reach a consensus. These applications primarily focus on communication among agents and do not depend on the external environment.

% This environment defines the rules, interfaces, and boundaries through which agents communicate and perform their tasks. 
% 2. Why the Environment concept exists in the multi-agent work / Why the environment is important in multi-agent?
% The Environment provides a fundamental structured medium through which agents can interact and communicate. 
% 3. Basic properties: What is state and space and Detail categorization of the basic properties
% There are two basic properties of the environment: State and Space. 

% \paragraph{State:} The State refers to whether the condition or status of the multi-agents system would change over time. The State can be categorized into three types: None, Static or Dynamic. None means that there is no specific environment setting. The Static means that the environment remains constant or unchanged over time, while the dynamic means that the environment can change or evolve over time. 

% \paragraph{Space:} The Space represents the type of environment or platform in which the agents operate and the space can be categorized into three types: Text-only, Sandbox, and Physical. Text-only means that the interactions and activities of the agents are conducted entirely through written text. In such a space, the agents’ capabilities are primarily focused on understanding, generating, and responding to textual information. The Sandbox refers to a simulated or virtual environment where agents can interact more freely and experiment with various actions and strategies. Such space allows for more complex interactions than a text-only space. The Physical is a real-world environment where agents interact with physical entities and obey real-world physics and constraints. In physical space, agents normally need to take actions that can have direct physical outcomes. 
% 4. How to implement this concept in large language model based multi-agents work (Usually)

\subsection{Agents Profiling}
% What is (Definition, suitable for which kind of scenario) / 
% Why need (insights from real world, pros,cons compared to other methods) / (potential limitations) / 
% How to implement in LLM /

In LLM-MA systems, agents are defined by their traits, actions, and skills, which are tailored to meet specific goals. Across various systems, agents assume distinct roles, each with comprehensive descriptions encompassing characteristics, capabilities, behaviors, and constraints. For instance, in gaming environments, agents might be profiled as players with varying roles and skills, each contributing differently to the game's objectives. In software development, agents could take on the roles of product managers and engineers, each with responsibilities and expertise that guide the development process. Similarly, in a debating platform, agents might be designated as proponents, opponents, or judges, each with unique functions and strategies to fulfill their roles effectively. These profiles are crucial for defining the agents' interactions and effectiveness within their respective environments. Table \ref{summary} lists the agent \emph{\textbf{Profiles}} in recent LLM-MA works.

%% 1.What is Agent Profiling
%For Agent Profiling, this concept involves identifying and comprehending the traits, actions, and skills of each agent. 
%% 2. Why the Agent Profiling concept exists in the multi-agent work / Why important?
%Agent Profiling is fundamental in multi-agent systems to help assign a unique role to each agent and ensure the collective capabilities of the system are fully leveraged. 
%% 3. Basic properties: What is state and space and Detail categorization of the basic properties
%There are two basic properties of Agent Profiling: Agent Profiling Methods and Agent Profiles. 

%\paragraph{Agent Profiling Methods:} The Agent Profiling Methods refer to the methods to build the agent profiles. The  Agent Profiling Methods can be 
Regarding the \emph{\textbf{Agent Profiling Methods}}, we 
categorized them into three types: \emph{Pre-defined, Model-Generated}, and \emph{Data-Derived}.
%From existing datasets. 
In the \emph{\textbf{Pre-defined}} cases,  agent profiles are explicitly defined by the system designers. 
The \emph{\textbf{Model-Generated}} method creates agent profiles by models, e.g., large language models. The \emph{\textbf{Data-Derived}} method involves constructing agent profiles based on pre-existing datasets.
%means that the agent profiles are created by existing datasets. 

%\paragraph{Agent Profiles:} 
% What is
%The Agent Profiles refer to the comprehensive descriptions of the roles (characteristics, capabilities, behaviors, constraints, etc) of individual agents within the system. The setting of the agent profiles depends on the practical problems.
% Why need
%Well-designed agent profiles can help each agent play its roles better and ensure that each agent is utilized according to its strengths, thus leading to effective coordination and task distribution.

% 4. How to implement this concept in large language model based multi-agents work (Usually)


% \subsection{Agent Management}

% % Basic concept
% % 1. What is (definition)
% % 2.  Why need (importance to multi-agent system, insights from real world)

% % Details concept
% % 1. What is (Definition, suitable for which kind of scenario)
% % 2. Why need (insights from real world, pros,cons compared to other methods) / (potential limitations) 
% % 3. How to implement in LLM

% % What is (Definition) 
% Agent Management refers to the strategies used to oversee or optimize the functioning of individual agents within the system. It involves how agents are created, maintained, and potentially evolved or decommissioned over time. 
% % Why need (importance to a multi-agent system; insights from real world)
% % Agent management ensures that the collective capabilities of the agents are fully utilized over time, and the multi-agent system as a whole can adapt and evolve in response to new challenges.

% % Basic categories: What is and Why need
% Agent Management can be categorized into three types based on the previous work: Static, Dynamic Generation. The Static means that agents are created with a fixed set of capabilities and roles that do not change over time. Static agent management is suitable for the case that agents continue to perform their designated functions reliably without the need for updates or modifications. Dynamic Generation involves creating new agents on-the-fly during the system’s operation and this management method allows the system to scale and adapt by introducing agents tailored to current needs. 
% Dynamic Refinement means that existing agents are capable of learning and improving over time based on their interactions and experiences. This method is ideal for environments that are dynamic and unpredictable. 



\subsection{Agents Communication} \vspace{-0.05cm}
% Basic concept
% 1. What is (definition)
%Agent Communication refers to how agents interact with each other in the multi-agent system to achieve individual or collective goals.
% 2.  Why need (importance to multi-agent system, insights from real world)
%Agent Communication is crucial for the multi-agent system because it is the backbone that supports collective intelligence.
% Details concept
% 1. What is (Definition, suitable for which kind of scenario)
% 2. Why need (insights from real world, pros,cons compared to other methods) / (potential limitations) 
% 3. How to implement in LLM
%There are two basic properties of Agent Communication: Communication Paradigms and Communication Structure.

The communication between agents in LLM-MA systems is the critical infrastructure supporting collective intelligence. We dissect agent communication from three perspectives: 1) \emph{Communication Paradigms}:   the styles and methods of interaction between agents; 2) \emph{Communication Structure}: the organization and architecture of communication networks within the multi-agent system; and 3) \emph{Communication Content}  exchanged between agents.

\paragraph{Communication Paradigms:} %The Communication Paradigms refers to the communication style of interaction between agents. The Communication Paradigms Methods can be categorized into three types: 
Current LLM-MA systems mainly take three paradigms for communication: 
\emph{Cooperative}, \emph{Debate}, and \emph{Competitive}.  \emph{\textbf{Cooperative}}  agents work together towards a shared goal or objectives, typically exchanging information to enhance a collective solution. The \emph{\textbf{Debate}} paradigm is employed when agents engage in argumentative interactions, presenting and defending their own viewpoints or solutions, and critiquing those of others. This  paradigm is ideal for reaching a consensus or a more refined solution. \emph{\textbf{Competitive}} agents work towards their own goals that might be in conflict with the goals of other agents. 


\begin{figure}[t]
\centering
\includegraphics[width= 0.44\textwidth]{img/Agent_communication.png}
\caption{The Agent Communication Structure.} 
\label{fig: Agent_communication} 
\end{figure}

\paragraph{Communication Structure:} %The Communication Structure refers to the way communication networks are organized within the multi-agent system. The Communication Structure can be categorized into three types: Layered, Decentralized, and Centralized. The details of these three communication structures are demonstrated in Fig.~\ref{fig: Agent_communication}. 
Fig.~\ref{fig: Agent_communication} shows {four} typical communication structures in LLM-MA systems. 
\emph{\textbf{Layered}} communication is structured hierarchically, with agents at each level having distinct roles and primarily interacting within their layer or with adjacent layers. \cite{liu2023dynamic} introduces a framework called Dynamic LLM-Agent Network (DyLAN), which organizes agents in a multi-layered feed-forward network. This setup facilitates dynamic interactions, incorporating features like inference-time agent selection and an early-stopping mechanism, which collectively enhance the efficiency of cooperation among agents. 
\emph{\textbf{Decentralized}} communication operates on a peer-to-peer network, where agents directly communicate with each other, a structure commonly employed in world simulation applications.
\emph{\textbf{Centralized}} communication involves a central agent or a group of central agents coordinating the system's communication, with other agents primarily interacting through this central node.
{\emph{\textbf{Shared Message Pool}} is proposed by MetaGPT~\cite{hong2023metagpt} to improve the communication efficiency. This communication structure maintains a shared message pool where agents publish messages and subscribe to relevant messages based on their profiles, thereby boosting communication efficiency.}
% examples
%means that the communication is organized in hierarchical layers, where agents in each layer have specific roles and communicate predominantly within that layer or with the adjacent layers. The Decentralized means that the communication is a peer-to-peer network and agents communicate directly with each other. This communication structure is widely used in world simulation applications. Centralized means that in this communication, there is a central agent or a set of central agents that coordinate communication within the system while other agents communicate primarily through this central node. 

\paragraph{Communication Content:} %The Communication Content refers to the communication information between agents. In LLM-based Multi-Agents systems, the type of communication information would be text. The concrete content between different agents would be diverse and 
In LLM-MA systems, the Communication Content typically takes the form of text. The specific content varies widely and  depends on the particular application. For example, in software development, agents may communicate with each other about code segments. In simulations of games like Werewolf,  agents might discuss their analyses, suspicions, or strategies.  

% Code, Game.. 


\subsection{Agents Capabilities Acquisition}
% How LLM-based Multi-Agents Learn during communication
% What learned: 1. Experience => Memory  2. On-policy: Temammates dynamically
% How the learned help 1. Self-revolution (Prompt?) Plan
% ProAgent: Building Proactive Cooperative AI with Large Language Models
The Agents Capabilities Acquisition is a crucial process in LLM-MA, enabling agents to learn and evolve dynamically. In this context, there are two fundamental concepts: the types of \emph{\textbf{feedback}} from which agents should learn to enhance their capabilities,  and the \emph{\textbf{strategies}} for agents to \emph{\textbf{adjust}} themselves to effectively solve complex problems. 
%\emph{\textbf{designing}} inter-agents modeling 
% to effectively solve complex problems.
%Based on previous introduced basic elements of LLM-based Multi-Agents, a LLM-based Multi-Agents system can be built. The Agents Capabilities Acquisition helps the agents in the multi-agent system learn and revolution dynamically, empowering the capabilities of the multi-agent system. For Capabilities Acquisition, two important basic concepts are Feedback and Agents Adjustment to Complex Problems. 
% We also summarize the advanced knowledge sharing strategy for Multi-Agents systems.
\paragraph{Feedback:} Feedback involves the critical information that agents receive about the outcome of their actions, helping the agents learn the potential impact of their actions and adapt to complex and dynamic problems. In most studies, the format of feedback provided to agents is textual. Based on the sources from which agents receive this feedback, it can be categorized into four types. \textbf{1) Feedback from Environment}, e.g.,  from either real world environments or   virtual environments~\cite{wang2023survey}. It is prevalent in most LLM-MA for problem-solving scenarios, including Software Development (agents obtain feedback from Code Interpreter), and Embodied multi-agents systems (robots obtain feedback from real-world or Simulated environments). \textbf{2) Feedback from Agents Interactions}  means that the feedback comes from the judgement of other agents or from agents communications. It is common in problem-solving scenarios like science debates, where agents learn to  critically evaluate and refine the conclusions through communications. In  world simulation scenarios such as Game Simulation,  agents learn to refine strategies based on previous interactions between other agents. \textbf{3)  Human Feedback} comes directly from humans and is crucial for aligning the multi-agent system with human values and preferences.   This kind of feedback is widely used in most ``Human-in-the-loop'' applications~\cite{wang2021putting}. Last \textbf{4) None}.  In some cases, there is no feedback provided to  the agents. This  often happens for world simulation works focused on analyzing simulated results rather than the planning capabilities of agents. In such scenarios, like propagation simulation, the emphasis is on result analysis, and hence, feedback is not a component of the system. 

\paragraph{Agents Adjustment to Complex Problems:} 
To enhance their capabilities, agents in LLM-MA systems can adapt through three main solutions.
%Multi-Agents systems are usually used in complex problem scenarios, most multi-agents systems need to dynamically adjust themselves to rapidly adapt to complex problems and perform better in the future. The adjustments to complex problems can be categorized into three types: Memory, Self-Evolution and Dynamic Generation. 
\textbf{1) Memory.} Most LLM-MA systems leverage a memory module for agents to adjust their behavior. Agents store information from previous interactions and feedback in their memory. When performing actions, they can retrieve relevant, valuable memories, particularly those containing successful actions for similar past goals, as highlighted in~\cite{wang2023survey}. This process aids in enhancing their current actions.  
\textbf{2) Self-Evolution.}  Instead of  {only relying on the historical records to decide subsequent actions as seen in  \emph{Memory}-based solutions, agents can dynamically self-evolve by modifying themselves such as altering their initial goals and planning strategies, and training themselves based on feedback or communication logs. \cite{nascimento2023self} proposes a self-control loop process to allow each agent in the multi-agents systems to be self-managed and self-adaptive to dynamic environments, thereby improving the cooperation efficiency of multiple agents. {
% In the LLM-MA, when each agent effectively models the goals, plans, and beliefs of other agents, it can help facilitate mutual understanding, thereby enhancing collaborative planning capabilities.
\cite{zhang2023proagent} introduces ProAgent which anticipates teammates' decisions and dynamically adjusts each agent's strategies based on the communication logs between agents, facilitating mutual understanding and improving collaborative planning capability.} \cite{wang2023adapting} discusses a  \emph{Learning through Communication} (LTC) paradigm, using the communication logs of multi-agents to generate datasets to train or fine-tune LLMs. LTC enables continuous adaptation and improvement of agents through interaction with their environments and other agents, breaking the limits of in-context learning or supervised fine-tuning, which don't fully utilize the feedback received during interactions with the environment and external tools for continuous training. Self-Evolution enables agents' autonomous adjustment in their profiles or goals, rather than just learning from historical interactions.} 
\textbf{3) Dynamic Generation.} In some scenarios,   the system can generate new agents on-the-fly during its operation~\cite{chen2023autoagents,chen2023agentverse}. This capability enables the system to scale and adapt effectively, as it can introduce agents that are specifically designed to address current needs and challenges. 

{With the scaling up LLM-MA with  a larger  number of agents,  the escalating complexity of managing various kinds of agents has been a critical problem. \emph{Agents Orchestration} emerged as a pivotal challenge and began to gain attention in ~\cite{multiorches1,multiorches2}. We will further discuss this topic in Section \ref{scale_up}.}
% The Self-Evolution method emphasizes on the self profiles or goals adjustment of the agents themselves instead of learning from history interactions like the Memory method. The Dynamic Generation involves creating new agents on-the-fly during the system’s operation and this management method allows the system to scale and adapt by introducing agents tailored to current needs.

% \paragraph{Knowledge Sharing:} The Feedback and Adjustment to Complex Problems are two basic concepts for each agent in the multi-agents system. For the synergy between multiple agents, MetaGPT~\cite{hong2023metagpt} proposes the Knowledge Sharing strategy to main a shared common environment and broadcast the learned knowledge between agents. By doing this, agents can exchange their learned knowledge effectively and enhance the collective learning and performance of the system.


% Relation between Agents and Environment
% 

% Knowledge sharing / partially sharing / non-sharing

% Information Mediums
% broadcasting

% \subsection{Training Scheme}


% Multi-Modal Agents Framework
%\paragraph{Multi-Agents Orchestration and Efficiency:} 
% \paragraph{Agents Orchestration:}

% Multi-agents orchestration and efficiency have been discussed to enhance the ability of agents to work together effectively and adapt to complex environments~\cite{multiorches1,multiorches2}. 


% \paragraph{Inter-Agents Modeling:}
%\subsection{Multi-Agents Orchestration and Efficiency}
%\label{Sec5:Orchestration}
% What and why efficient communiction is important
%In previous taxonomy of LLM-based Multi-Agents systems, we introduce that multiple agents can have different profiles, communicate with each other, and acquire capabilities from different feedback. These are basics for the LLM Multi-Agents systems but may not maximize the capabilities of the system. Hence, 
% To harness and maximize the potential benefits of the collective intelligence in LLM-MA systems, \xz{promoting collaborative learning is essential. This concept is widely discussed in previous Multi-Agents systems and Reinforcement Learning based Multi-Agents systems,} 





\begin{table*}[h!]
\centering
\scalebox{0.59}{  
\begin{tabular}{c|cc|c|c|cc|cc|cc}
\toprule
 & &    &    &  {\color[HTML]{BF9000} } &
  \multicolumn{2}{c|}{{\color[HTML]{4472C4} \textbf{ Agents   Profiling }}} &
  \multicolumn{2}{c|}{{\color[HTML]{548135} \textbf{\begin{tabular}[c]{@{}c@{}}Agents \\ Communication\end{tabular}}}} &
  \multicolumn{2}{c}{{\color[HTML]{FF00FF} \textbf{Agents Capabilities Acquisition}}} \\ \cline{6-11} 
\multirow{-2}{*}{\textbf{Motivation}} &
  \multicolumn{2}{c|}{\multirow{-2}{*}{\textbf{Research Domain \& Goals}}} &
  \multirow{-2}{*}{\textbf{Work}} &
  \multirow{-4}{*}{{\color[HTML]{BF9000} \textbf{\begin{tabular}[c]{@{}c@{}}Agents-Env.\\ Interface  \end{tabular}}}} &
  %\\  ( Sandbox /  \\ Physical /\\ None )\end{tabular}}}} &
  \multicolumn{1}{c|}{{\color[HTML]{4472C4} \textbf{ \begin{tabular}[c]{@{}c@{}}Profiling\\ methods\end{tabular}}   }} &  %( Pre-defined / \\ Model-Generated / \\ Data-Derived )\end{tabular}}}}   &
  {\color[HTML]{4472C4} \textbf{\begin{tabular}[c]{@{}c@{}}Profiles\\ (examples)\end{tabular}}} &
  \multicolumn{1}{c|}{{\color[HTML]{548135} \textbf{Paradigms}}} &
%  \textbf{\begin{tabular}[c]{@{}c@{}}  Paradigms\\ ( Cooperative / \\ Debate / \\ Competitive )\end{tabular}}}} &
  {\color[HTML]{548135}  \textbf{Structure}} &
%  \textbf{\begin{tabular}[c]{@{}c@{}}Structure\\ ( Layered /\\ Decentralized / \\ Centralized)\end{tabular}}} &
  \multicolumn{1}{c|}{{\color[HTML]{FF00FF} \textbf{Feedback from}}} &
  %\textbf{\begin{tabular}[c]{@{}c@{}}Feedback from\\ (Environment /\\ Agents Interactions /\\ Human)\end{tabular}}}} &
  {\color[HTML]{FF00FF} \textbf{\begin{tabular}[c]{@{}c@{}}Agents \\ Adjustment \end{tabular}}}     \vspace{1pt} \\ \hline 
  %\\ (Memory update from experience /\\ Self-Evolution from experience /\\ Dynamic Generation)\end{tabular}}} \\ \hline
 &
  \multicolumn{2}{c|}{} &
 {\scriptsize  \cite{qian2023communicative}} &
  Sandbox &
  \multicolumn{1}{c|}{\begin{tabular}[c]{@{}c@{}}Pre-defined, \\ Model-Generated\end{tabular}} &
  \begin{tabular}[c]{@{}c@{}}CTO,\\ programmer\end{tabular} &
  \multicolumn{1}{c|}{Cooperative} &
  Layered &  
  \multicolumn{1}{c|}{   \begin{tabular}[c]{@{}c@{}}Environment, \\ Agent interaction,\\ Human\end{tabular}} &
  \begin{tabular}[c]{@{}c@{}}Memory,\\ Self-Evolution\end{tabular} \\ \cline{4-11} 
 &
  \multicolumn{2}{c|}{Software development} &
  { \scriptsize \cite{hong2023metagpt}} &
  Sandbox &
  \multicolumn{1}{c|}{Pre-defined} &
  \begin{tabular}[c]{@{}c@{}}Product Manager, \\ Engineer\end{tabular} &
  \multicolumn{1}{c|}{Cooperative} &
   \begin{tabular}[c]{@{}c@{}}Layered, \\ Shared Message Pool\end{tabular} &
  \multicolumn{1}{c|}{ \begin{tabular}[c]{@{}c@{}}Environment, \\ Agent interaction,\\ Human\end{tabular}} &
  \begin{tabular}[c]{@{}c@{}}Memory,\\ Self-Evolution\end{tabular} \\ \cline{4-11} 
 &
  \multicolumn{2}{c|}{ } &
  { \scriptsize \cite{dong2023selfcollaboration}} &
  Sandbox &
  \multicolumn{1}{c|}{\begin{tabular}[c]{@{}c@{}}Pre-defined, \\ Model-Generated\end{tabular}} &
  \begin{tabular}[c]{@{}c@{}}Analyst,\\ coder\end{tabular} &
  \multicolumn{1}{c|}{Cooperative} &
  Layered &
  \multicolumn{1}{c|}{\begin{tabular}[c]{@{}c@{}}Environment,\\  Agent interaction\end{tabular}} &
  \begin{tabular}[c]{@{}c@{}}Memory,\\ Self-Evolution\end{tabular} \\ \cline{2-11} 
 &
  \multicolumn{1}{c|}{} & \begin{tabular}[c]{@{}c@{}}Multi-robot  \\ planning \end{tabular} &
 % Scalable multi-robot planning &
 { \hspace{-10pt}  \scriptsize  \cite{chen2023scalable} \hspace{-10pt} } &
  \begin{tabular}[c]{@{}c@{}}Sandbox, \\ Physical\end{tabular} &
  \multicolumn{1}{c|}{Pre-defined} &
  Robots &
  \multicolumn{1}{c|}{Cooperative} &
  \begin{tabular}[c]{@{}c@{}}Centralized,\\ Decentralized\end{tabular} &
  \multicolumn{1}{c|}{\begin{tabular}[c]{@{}c@{}}Environment,\\  Agent interaction\end{tabular}} &
  Memory \\ \cline{3-11} 
 &
 % \multicolumn{1}{c|}{\multirow{-3}{*}{Embodied Agents}} &
    \multicolumn{1}{c|}{\begin{tabular}[c]{@{}c@{}}Embodied\\ Agents \end{tabular} } &
   \begin{tabular}[c]{@{}c@{}}Multi-robot  \\ collaboration \end{tabular} &
%  Multi-robot collaboration &
 { \scriptsize \cite{mandi2023roco}} &
  \begin{tabular}[c]{@{}c@{}}Sandbox, \\ Physical\end{tabular} &
  \multicolumn{1}{c|}{Pre-defined} &
  Robots &
  \multicolumn{1}{c|}{Cooperative} &
  Decentralized &
  \multicolumn{1}{c|}{\begin{tabular}[c]{@{}c@{}}Environment,\\ Agent interaction\end{tabular}} &
  Memory \\ \cline{3-11} 
 &
  \multicolumn{1}{c|}{} &
 \begin{tabular}[c]{@{}c@{}}Multi-Agents  \\ cooperation \end{tabular} &
%  Multi-Agents cooperation &
 {\hspace{-15pt} \scriptsize \cite{zhang2023building} \hspace{-15pt}} &
  Sandbox &
  \multicolumn{1}{c|}{Pre-defined} &
  Robots &
  \multicolumn{1}{c|}{Cooperative} &
  Decentralized &
  \multicolumn{1}{c|}{\begin{tabular}[c]{@{}c@{}}Environment,\\ Agent interaction\end{tabular}} &
  Memory \\ \cline{2-11} 
 \multirow{-5}{*}{\begin{tabular}[c]{@{}c@{}}Problem \\ Solving\end{tabular}}  &
%  \multicolumn{1}{c|}{Science Experiments} &
   \multicolumn{1}{c|}{\begin{tabular}[c]{@{}c@{}} Science \\ Experiments \end{tabular}} &
   \begin{tabular}[c]{@{}c@{}} Optimization \\ of MOF \end{tabular} &
%  Optimization of MOF. &
 {\hspace{-10pt} \scriptsize  \cite{chatgptforchem} \hspace{-10pt} } &
  Physical &
  \multicolumn{1}{c|}{Pre-defined} &
  \begin{tabular}[c]{@{}c@{}}Strategy planers,\\ literature \\ collector, coder\end{tabular} &
  \multicolumn{1}{c|}{Cooperative} &
  Centralized &
  \multicolumn{1}{c|}{\begin{tabular}[c]{@{}c@{}}Environment,\\ Human\end{tabular}} &
  Memory \\ \cline{2-11} 
 &
  \multicolumn{1}{c|}{} &
   \begin{tabular}[c]{@{}c@{}} Improving \\ Factuality \end{tabular} &
%  Improving Factuality &
 { \scriptsize \cite{du2023improving}} &
  None &
  \multicolumn{1}{c|}{Pre-defined} &
  Agents &
  \multicolumn{1}{c|}{Debate} &
  Decentralized &
  \multicolumn{1}{c|}{Agent interaction} &
  Memory \\ \cline{3-11} 
 &
    \multicolumn{1}{c|}{\begin{tabular}[c]{@{}c@{}}Science \\ Debate\end{tabular}} &
   \begin{tabular}[c]{@{}c@{}}Examining,\\ Inter-Consistency\end{tabular} &
%  Examining Inter-Consistency &
 {\hspace{-15pt} \scriptsize \cite{xiong2023examining} \hspace{-10pt}} &
  None &
  \multicolumn{1}{c|}{Pre-defined} &
  \begin{tabular}[c]{@{}c@{}}Proponent,\\ Opponent,\\ Judge\end{tabular} &
  \multicolumn{1}{c|}{Debate} &
  \begin{tabular}[c]{@{}c@{}}Centralized,\\ Decentralized\end{tabular} &
  \multicolumn{1}{c|}{Agent interaction} &
  Memory \\ \cline{3-11} 
 &
  \multicolumn{1}{c|}{} &
  \begin{tabular}[c]{@{}c@{}}Evaluators \\ for debates \end{tabular} &
%  Evaluators for debates &
 { \hspace{-10pt} \scriptsize \cite{chan2023chateval}} &
  None &
  \multicolumn{1}{c|}{Pre-defined} &
  Agents &
  \multicolumn{1}{c|}{Debate} &
  \begin{tabular}[c]{@{}c@{}}Centralized,\\ Decentralized\end{tabular} &
  \multicolumn{1}{c|}{Agent interaction} &
  Memory \\ \cline{3-11} 
&
  %\multicolumn{1}{c|}{\multirow{-4}{*}{Science Debate}} &
    \multicolumn{1}{c|}{} &
    \begin{tabular}[c]{@{}c@{}}Multi-Agents \\ for Medication\end{tabular} &
%  Multi-Agents for Medication &
 {\hspace{-10pt}  \scriptsize \cite{tang2023medagents} \hspace{-10pt} } &
  None &
  \multicolumn{1}{c|}{Pre-defined} &
  \begin{tabular}[c]{@{}c@{}}Cardiology,\\ Surgery\end{tabular} &
  \multicolumn{1}{c|}{\begin{tabular}[c]{@{}c@{}}Debate,\\ Cooperative\end{tabular}} &
  \begin{tabular}[c]{@{}c@{}}Centralized,\\ Decentralized\end{tabular} &
  \multicolumn{1}{c|}{Agent interaction} &
  Memory \\ \hline
 &
  \multicolumn{1}{c|}{} &
    \begin{tabular}[c]{@{}c@{}} Modest Community \\ (25 persons)\end{tabular} &
%  Modest Community (25 persons) &
 {\hspace{-10pt}  \scriptsize \cite{park2023generative} \hspace{-10pt} } &
  Sandbox &
  \multicolumn{1}{c|}{Model-generated} &
  \begin{tabular}[c]{@{}c@{}}Pharmacy,\\ shopkeeper\end{tabular} &
  \multicolumn{1}{c|}{-} &
  - &
  \multicolumn{1}{c|}{\begin{tabular}[c]{@{}c@{}}Environment,\\ Agent interaction\end{tabular}} &
  Memory \\ \cline{3-11} 
 &
  \multicolumn{1}{c|}{} &
  \begin{tabular}[c]{@{}c@{}}  Online community \\ (1000 persons)\end{tabular} &
%  Online community (1000 personas) &
 {\hspace{-10pt}  \scriptsize \cite{park2022social} \hspace{-10pt} } &
  None &
  \multicolumn{1}{c|}{\begin{tabular}[c]{@{}c@{}}Pre-defined,\\ Model-generated\end{tabular}} &
  \begin{tabular}[c]{@{}c@{}}Camping,\\  fishing\end{tabular} &
  \multicolumn{1}{c|}{-} &
  - &
  \multicolumn{1}{c|}{Agent interaction} &
 \begin{tabular}[c]{@{}c@{}}Dynamic\\  Generation\end{tabular}    \\ \cline{3-11} 
 &
  \multicolumn{1}{c|}{Society} &
  Emotion propagation &
 {\hspace{-10pt}  \scriptsize \cite{gao2023s} \hspace{-10pt} } &
  None &
  \multicolumn{1}{c|}{\begin{tabular}[c]{@{}c@{}}Pre-defined,\\ Model-generated\end{tabular}} &
  \begin{tabular}[c]{@{}c@{}}Real-world \\ user \end{tabular} &
  \multicolumn{1}{c|}{-} &
  - &
  \multicolumn{1}{c|}{Agent interaction} &
  Memory \\ \cline{3-11} 
 &
  \multicolumn{1}{c|}{} &
    \begin{tabular}[c]{@{}c@{}}  Real-time \\ social interactions \end{tabular} &
%  Real-time social interactions &
 {\hspace{-15pt} \scriptsize \cite{kaiya2023lyfe} \hspace{-15pt}} &
  Sandbox &
  \multicolumn{1}{c|}{Pre-defined} &
  \begin{tabular}[c]{@{}c@{}}Real-world \\ user \end{tabular} &
  \multicolumn{1}{c|}{-} &
  - &
  \multicolumn{1}{c|}{\begin{tabular}[c]{@{}c@{}}Environment,\\ Agent interaction\end{tabular}} &
  Memory \\ \cline{3-11} 
 &
  \multicolumn{1}{c|}{ } &
  Opinion dynamics &
 {\hspace{-10pt}  \scriptsize \cite{li2023quantifying}} &
  None &
  \multicolumn{1}{c|}{Pre-defined} &
  \begin{tabular}[c]{@{}c@{}}NIN, NINL,\\ NIL\end{tabular} &
  \multicolumn{1}{c|}{-} &
  - &
  \multicolumn{1}{c|}{Agent interaction} &
  Memory \\ \cline{2-11} 
 &
  \multicolumn{1}{c|}{} &
  WereWolf &
  \begin{tabular}[c]{@{}c@{}} {\scriptsize \cite{xu2023exploring} }\\ {\scriptsize  \cite{xu2023language} }\end{tabular} &
  Sandbox &
  \multicolumn{1}{c|}{Pre-defined} &
  \begin{tabular}[c]{@{}c@{}}Seer,\\ werewolf,\\ villager\end{tabular} &
  \multicolumn{1}{c|}{\begin{tabular}[c]{@{}c@{}}Cooperative, \\ Debate,\\ Competitive\end{tabular}} &
  Decentralized &
  \multicolumn{1}{c|}{\begin{tabular}[c]{@{}c@{}}Environment,\\ Agent interaction\end{tabular}} &
  Memory \\ \cline{3-11} 
 &
  \multicolumn{1}{c|}{Gaming} &
  Avalon &
  \begin{tabular}[c]{@{}c@{}} {\hspace{-10pt} \scriptsize \cite{light2023avalonbench} \hspace{-10pt} } \\ {\hspace{-10pt} \scriptsize \cite{wang2023avalon} \hspace{-10pt} } \end{tabular} &
  Sandbox &
  \multicolumn{1}{c|}{Pre-defined} &
  \begin{tabular}[c]{@{}c@{}}Servant,\\ Merlin,\\  Assassin\end{tabular} &
  \multicolumn{1}{c|}{\begin{tabular}[c]{@{}c@{}}Cooperative, \\ Debate,\\ Competitive\end{tabular}} &
  Decentralized &
  \multicolumn{1}{c|}{\begin{tabular}[c]{@{}c@{}}Environment,\\ Agent interaction\end{tabular}} &
  Memory \\ \cline{3-11} 
 &
  \multicolumn{1}{c|}{} &
  Welfare Diplomacy &
 {\hspace{-10pt}  \scriptsize \cite{mukobi2023welfare} \hspace{-10pt}} &
  Sandbox &
  \multicolumn{1}{c|}{Pre-defined} &
  Countries &
  \multicolumn{1}{c|}{\begin{tabular}[c]{@{}c@{}}Cooperative,\\ Competitive\end{tabular}} &
  Decentralized &
  \multicolumn{1}{c|}{\begin{tabular}[c]{@{}c@{}}Environment,\\ Agent interaction\end{tabular}} &
  Memory \\ \cline{2-11} 
 &
  \multicolumn{1}{c|}{} &
    \begin{tabular}[c]{@{}c@{}}  Human behavior \\ Simulation\end{tabular} &
%  Human behavior Simulation &
{\hspace{-10pt}  \scriptsize  \cite{aher2023using} } &
  Sandbox &
  \multicolumn{1}{c|}{Pre-defined} &
  Humans &
  \multicolumn{1}{c|}{-} &
  - &
  \multicolumn{1}{c|}{Agent interaction} &
  Memory \\ \cline{3-11} 
\multirow{-18}{*}{\begin{tabular}[c]{@{}c@{}}World \\ \hspace{-10pt} Simulation \hspace{-10pt} \end{tabular}} 
 &
  \multicolumn{1}{c|}{\multirow{-2}{*}{Psychology}} &
  \begin{tabular}[c]{@{}c@{}}Collaboration\\ Exploring \end{tabular} &
%  Collaboration Exploring &
  {\hspace{-10pt} \scriptsize \cite{zhang2023exploring} \hspace{-10pt} } &
  None &
  \multicolumn{1}{c|}{Pre-defined} &
  Agents &
  \multicolumn{1}{c|}{\begin{tabular}[c]{@{}c@{}}Cooperative,\\ Debate\end{tabular}} &
  Decentralized &
  \multicolumn{1}{c|}{Agent interaction} &
  Memory \\ \cline{2-11} 
 &
  \multicolumn{1}{c|}{} &
    \begin{tabular}[c]{@{}c@{}}Macroeconomic \\ simulation \end{tabular} &
%  Macroeconomic simulation &
  {\hspace{-10pt}  \scriptsize \cite{li2023large}} &
  None &
  \multicolumn{1}{c|}{\begin{tabular}[c]{@{}c@{}}Pre-defined,\\ Model-generated\end{tabular}} &
  Labor &
  \multicolumn{1}{c|}{Cooperative} &
  Decentralized &
  \multicolumn{1}{c|}{Agent interaction} &
  Memory \\ \cline{3-11} 
 &
  \multicolumn{1}{c|}{Economy} &
  \begin{tabular}[c]{@{}c@{}}Information \\ Marketplaces \end{tabular} &
 % Information Marketplaces &
 {\hspace{-10pt}  \scriptsize \cite{anonymous2023rethinking} \hspace{-10pt} } &
  Sandbox &
  \multicolumn{1}{c|}{\begin{tabular}[c]{@{}c@{}}Pre-defined, \\ Data-Derived\end{tabular}} &
  Buyer &
  \multicolumn{1}{c|}{\begin{tabular}[c]{@{}c@{}}Cooperative, \\ Competitive\end{tabular}} &
  Decentralized &
  \multicolumn{1}{c|}{\begin{tabular}[c]{@{}c@{}}Environment,\\ Agent interaction\end{tabular}} &
  Memory \\ \cline{3-11} 
 &
  \multicolumn{1}{c|}{} &
    \begin{tabular}[c]{@{}c@{}}Improving \\ financial trading \end{tabular} &

%  Improving financial trading &
  {\hspace{-10pt}  \scriptsize \cite{li2023tradinggpt}} &
  Physical &
  \multicolumn{1}{c|}{Pre-defined} &
  Trader &
  \multicolumn{1}{c|}{Debate} &
  Decentralized &
  \multicolumn{1}{c|}{\begin{tabular}[c]{@{}c@{}}Environment,\\ Agent interaction\end{tabular}} &
  Memory \\ \cline{3-11} 
 &
  \multicolumn{1}{c|}{} &
  Economic theories &
  {\hspace{-10pt}  \scriptsize \cite{zhao2023competeai}} &
  Sandbox &
  \multicolumn{1}{c|}{\begin{tabular}[c]{@{}c@{}}Pre-defined, \\ Model-Generated\end{tabular}} &
  \begin{tabular}[c]{@{}c@{}}Restaurant,\\  Customer\end{tabular} &
  \multicolumn{1}{c|}{Competitive} &
  Decentralized &
  \multicolumn{1}{c|}{\begin{tabular}[c]{@{}c@{}}Environment,\\ Agent interaction\end{tabular}} &
  \begin{tabular}[c]{@{}c@{}}Memory,\\ Self-Evolution\end{tabular} \\ \cline{2-11} 
 &
  \multicolumn{1}{c|}{\multirow{2}{*}{ \begin{tabular}[c]{@{}c@{}} \hspace{-7pt} Recommender \hspace{-5pt}  \\ Systems \end{tabular}}} &
  \begin{tabular}[c]{@{}c@{}}Simulating \\ user behaviors\end{tabular} &
  %Simulating  user behaviors &
  {\hspace{-15pt} \scriptsize \cite{zhang2023generative} \hspace{-10pt}} &
  Sandbox &
  \multicolumn{1}{c|}{Data-Derived} &
  \begin{tabular}[c]{@{}c@{}}Users  from \\ MovieLens-1M\end{tabular} &
  \multicolumn{1}{c|}{-} &
  - &
  \multicolumn{1}{c|}{Environment} &
  Memory \\ \cline{3-11} 
 &
  \multicolumn{1}{c|}{ } &
    \begin{tabular}[c]{@{}c@{}}\hspace{-10pt} Simulating user-item \hspace{-10pt} \\ interactions \end{tabular} &
%  Simulating user-item interactions &
  { \hspace{-15pt} \scriptsize \cite{zhang2023agentcf} \hspace{-10pt}} &
  Sandbox &
  \multicolumn{1}{c|}{\begin{tabular}[c]{@{}c@{}}Pre-defined, \\ Data-Derived\end{tabular}} &
  \begin{tabular}[c]{@{}c@{}}User Agents\\ Item Agents\end{tabular} &
  \multicolumn{1}{c|}{Cooperative} &
  Decentralized &
  \multicolumn{1}{c|}{\begin{tabular}[c]{@{}c@{}}Environment,\\ Agent interaction\end{tabular}} &
  Memory \\ \cline{2-11} 
 &
  \multicolumn{1}{c|}{\multirow{2}{*}{ \begin{tabular}[c]{@{}c@{}} Policy \\ Making \end{tabular}}} &
  \begin{tabular}[c]{@{}c@{}}Public \\ Administration \end{tabular} &
%  Public Administration &
  {\hspace{-10pt}  \scriptsize \cite{xiao2023simulating}} &
  None &
  \multicolumn{1}{c|}{Pre-defined} &
  Residents &
  \multicolumn{1}{c|}{Cooperative} &
  Decentralized &
  \multicolumn{1}{c|}{Agent interaction} &
  Memory \\ \cline{3-11}  
 &
  \multicolumn{1}{c|}{} &
    War Simulation &
  {\hspace{-10pt}  \scriptsize \cite{hua2023war}} &
  None &
  \multicolumn{1}{c|}{Pre-defined} &
  Countries &
  \multicolumn{1}{c|}{Competitive} &
  Decentralized &
  \multicolumn{1}{c|}{Agent interaction} &
  Memory \\ \cline{2-11} 
  & \multicolumn{1}{c|}{\multirow{2}{*}{ \begin{tabular}[c]{@{}c@{}} Disease \end{tabular}}} &
   \begin{tabular}[c]{@{}c@{}}Human Behaviors  \\ to epidemics\end{tabular} &
 % Human Behaviors to epidemics &
 {\hspace{-6pt} \scriptsize \begin{tabular}[c]{@{}c@{}} [Ghaffarzadegan \\ \emph{et al.}, 2023] \end{tabular}}    &
% \cite{ghaffarzadegan2023generative}   &
  Sandbox &
  \multicolumn{1}{c|}{\begin{tabular}[c]{@{}c@{}}Pre-defined, \\ Model-Generated\end{tabular}} &
 \begin{tabular}[c]{@{}c@{}}Conformity   \\ traits \end{tabular} &
  \multicolumn{1}{c|}{Cooperative} &
  Decentralized &
  \multicolumn{1}{c|}{\begin{tabular}[c]{@{}c@{}}Environment,\\ Agent interaction\end{tabular}} &
  Memory  \\ \cline{3-11}  \vspace{1pt} 
  & \multicolumn{1}{c|}{} &
  Public health &
  { \scriptsize \begin{tabular}[c]{@{}c@{}} [Williams \\ \emph{et al.}, 2023] \end{tabular}} &
 % \cite{williams2023epidemic}} &
  Sandbox &
  \multicolumn{1}{c|}{\begin{tabular}[c]{@{}c@{}}Pre-defined, \\ Model-Generated\end{tabular}} &
  \begin{tabular}[c]{@{}c@{}} Adults aged  \\ 18 to 64\end{tabular} &
  \multicolumn{1}{c|}{Cooperative} &
  Decentralized &
  \multicolumn{1}{c|}{\begin{tabular}[c]{@{}c@{}}Environment,\\ Agent interaction\end{tabular}} &
  \begin{tabular}[c]{@{}c@{}}Memory,\\ Dynamic \\ Generation\end{tabular} \\ \bottomrule
\end{tabular}
}
\caption{Summary of the LLM-MA studies. We categorize current work according to their motivation, research domains and goals, and detail each work from different aspects regarding %four key LLM-based Multi-Agents elements: 
Agents-Environment Interface, Agents Profiling, Agents Communication and Agents Capability Acquisition. 
``-" denotes that a particular element is not specifically mentioned in this work.}
\label{summary}
\end{table*}

\section{Applications} 
\label{Sec5:Application}
% Practical applications (We can discuss each application from Agent specific Features, Agent organization, environment, ……. etc perspectives. Some applications may only emphasize some perspectives (not all perspectives). For example, there is no environment in multi-agent software development).  
% Overview Table (Taicheng)
% Please add the following required packages to your document preamble:
% \usepackage{multirow}

% Please add the following required packages to your document preamble:
% \usepackage{multirow}
% \usepackage[table,xcdraw]{xcolor}
% Beamer presentation requires \usepackage{colortbl} instead of \usepackage[table,xcdraw]{xcolor}
% Please add the following required packages to your document preamble:
% \usepackage{multirow}
% \usepackage[table,xcdraw]{xcolor}
% Beamer presentation requires \usepackage{colortbl} instead of \usepackage[table,xcdraw]{xcolor}

LLM-MA systems have been used in a wide range of applications. We summarize two kinds of applications in Table \ref{summary}: \textbf{Problem Solving} and  \textbf{World Simulation}. We   elaborate on these applications below. Note that this is a fast growing research field and new applications appear almost everyday. We maintain an  \href{https://github.com/taichengguo/LLM_MultiAgents_Survey_Papers}{open source repository} to report the latest work.  

\subsection{LLM-MA for Problem Solving}

% What is for problem solving and some simple examples
%One main stream application of LLM-MA systems is in specific problem-solving scenarios, such as Software Development, Embodied Agents, Science Experiments and Science Debate.
The main motivation of using LLM-MA for problem solving is to harness the collective capabilities of agents with specialized expertise. These agents, each acting as individuals, collaborate to address complex problems effectively,  such as software development, embodied agents, science experiments and science debate. These application examples are introduced next.
% General Methodology: How to achieve this from four taxonomy perspective
% In the context of LLM-MA for problem solving, the environment is typically Sandbox or Physical . Each agent assumes a distinct, specialized role, working cooperatively to solve the target task.
% More details as follows
%Subsequent sections delve into a detailed explanation of LLM-MA for problem solving applications.


\subsubsection{Software Development}
% Summary: 1.Why this problem needs LLM-MA  2. How LLM-MA do in this case (Schema from section 3)
Given that software development is a complex endeavor requiring the collaboration of various roles like product managers, programmers, and testers, LLM-MA systems are typically set to emulate these distinct roles and collaborate to address the intricate challenge. Following the waterfall or Standardized Operating Procedures (SOPs) workflow of the software development, the communication structure among agents is usually layered. Agents generally interact with the code interpreter, other agents or human to iteratively refine the generated code.
% Each work
\cite{li2023camel} first proposes a simple role-play agent framework, which utilizes the interplay of two roles to realize autonomous programming based on one-sentence user instruction. It provides insights into the ``cognitive" processes of communicative agents. \cite{dong2023selfcollaboration} makes LLMs work as distinct ``experts" for sub-tasks in software development, autonomously collaborating to generate code. Moreover, \cite{qian2023communicative} presents an end-to-end framework for software development, utilizing multiple agents for software development without incorporating advanced human teamwork experience. \cite{hong2023metagpt} first incorporates human workflow insights for more controlled and validated performance. It encodes SOPs into prompts to enhance structured coordination. \cite{huang2023agentcoder} delves deeper into multi-agent based programming by solving the problem of balancing code snippet generation with effective test case generation, execution, and optimization. 

% Some key limitations for LLM-MA for Software Development discussed in previous work are:
% There are some general limitations and challenges. Notably, observed limitations include inherent randomness in generated output, resulting in variability across software productions in different runs. Furthermore, the LLMs may manifest biases, deviating from problem-solving approaches typical of real programmers. This bias introduces code patterns not necessarily aligned with authentic programming paradigms. These challenges are further compounded by the hallucinatory tendencies inherent in large language models, underscoring the complexities associated with ensuring accurate and unbiased code generation in the realm of automatic programming.

\subsubsection{Embodied Agents}
% Summary: 1.Why this problem needs LLM-MA  2. How LLM-MA do in this case (Schema from section 3)
Most embodied agents applications inherently utilize multiple robots working together to perform complex real-world planning and manipulation tasks such as warehouse management with heterogeneous robot capabilities. Hence, LLM-MA can be used to model robots with different capabilities and cooperate with each other to solve real-world physical tasks.
% Each Work
\cite{dasgupta2023collaborating} first explores the potential to use LLM as an action planner for embedded agents. \cite{mandi2023roco} introduces RoCo, a novel approach for multi-robot collaboration that uses LLMs for high-level communication and low-level path planning. Each robotic arm is equipped with an LLM, cooperating with inverse kinematics and collision checking. Experimental results demonstrate the adaptability and success of RoCo in collaborative tasks. \cite{zhang2023building} presents CoELA, a Cooperative Embodied Language Agent, managing discussions and task planning in an LLM-MA setting. This challenging setting is featured with decentralized control, complex partial observation, costly communication, and multi-objective long-horizon tasks. \cite{chen2023scalable} investigates communication challenges in scenarios involving a large number of robots, as assigning each robot an LLM will be costly and unpractical due to the long context. The study compares four communication frameworks, centralized, decentralized, and two hybrid models, to evaluate their effectiveness in coordinating complex multi-agent tasks. \cite{yu2023conavgpt} proposes Co-NavGPT for multi-robot cooperative visual target navigation, integrating LLM as a global planner to assign frontier goals to each robot. \cite{chen2023multi} proposes an LLM-based consensus-seeking framework, which can be applied as a cooperative planner to a multi-robot aggregation task.


% The limitations 
% There are some general limitations. Current works for LLMs-based Embodied Agents mainly focus on high-level task planning, considering the inefficiency and poor performance of LLMs for low-level actions \cite{zhang2023building}. Besides, it is usually needed that the environment is already grounded in language \cite{nottingham2023embodied}, and multi-modal large models capable of both processing visual modalities effectively and generating natural language fluently \cite{huang2023language}\cite{driess2023palm} \cite{lu2022unified} would help. It is also mentioned that multiple rounds of discussions in a multi-agent system are inefficient where they might have difficulty reaching a consensus or repeat historical ideas and wrong actions \cite{chen2023scalable}.

\subsubsection{Science Experiments}
% ChatGPT team; Medical Team; 
% Motivation: Why use LLM-based Multi-Agents to solve this problem
Like multiple agents play as different specialists and cooperate to solve the Software Development and Embodied Agents problem, multiple agents can also be used to form a science team to conduct science experiments. One important difference from previous applications lies in the crucial role of human oversight,  due to the high expenses of the science experiments and the hallucination of the LLM agents. Human experts are at the center of these agents to process the information of agents and give feedback to the agents.
% How they use Multi-Agents; Conclusion
~\cite{chatgptforchem} utilizes multiple LLM-based agents, each focusing on specific tasks for the science experiments including strategy planning, literature search, coding, robotic operations, and labware design. All these agents interact with humans to work collaboratively to optimize the synthesis process of complex materials.

\subsubsection{Science Debate}
LLM-MA can be set for science debating scenarios, where agents debate with each other to enhance the collective reasoning capabilities in tasks such as Massive Multitask Language Understanding (MMLU)~\cite{hendrycks2020measuring}, Math problems~\cite{cobbe2021training}, and StrategyQA~\cite{geva2021did}. The main idea is that each agent initially offers its own analysis of a problem, which is then followed by a joint debating process. Through multiple rounds of debate, the agents converge on a single, consensus answer. 
~\cite{du2023improving} leverages the multi-agents debate process on a set of six different reasoning and factual accuracy tasks and demonstrates that LLM-MA debating can improve factuality. ~\cite{xiong2023examining} focuses on the commonsense reasoning tasks and formulates a three-stage debate to align with real-world scenarios including fair debate, mismatched debate, and roundtable debate. The paper also analyzes the inter-consistency between different LLMs and claims that debating can improve the inter-consistency.~\cite{tang2023medagents} also utilizes multiple LLM-based agents as distinct domain experts to do the collaborative discussion on a medical report to reach a consensus for medical diagnosis. 


\subsection{LLM-MA for World Simulation}

% What is for problem solving and some simple examples
Another mainstream application scenario of LLM-MA is the world simulation.  Research in this area is rapidly growing and spans a diverse range of fields including social sciences, gaming, psychology, economics, policy-making, etc.
% Why we need Multi-Agents in this case: LLM strong ability
The key reason for employing LLM-MA in world simulations lies in their exceptional role-playing abilities, which are crucial for realistically depicting various roles and viewpoints in a simulated world.
% General Methodology: How to achieve this from four taxonomy perspective
The environment of world simulation projects is usually crafted to reflect the specific scenario being simulated, with agents designed in various profiles to match this context. Unlike the problem solving systems that focus on agent cooperation, world simulation systems involve diverse methods of agent management and communication, reflecting the complexity and variety of real-world interactions.
% More details as follows
Next, we explore simulations conducted in diverse fields.
%The following sections provide in-depth information and analysis.



\subsubsection{Societal Simulation}
In societal simulation, LLM-MA models are used to simulate social behaviors, aiming to explore the potential social dynamics and propagation, test social science theories, and populate virtual spaces and communities with realistic social phenomena~\cite{park2023generative}.  Leveraging LLM's capabilities, agents with unique profiles engage in extensive communication, generating rich behavioral data for in-depth social science analysis. 

% Current Work
The scale of societal simulation has expanded over time, beginning with smaller, more intimate settings and progressing to larger, more intricate ones. 
Initial work by \cite{park2023generative} introduces generative agents within an interactive sandbox environment reminiscent of the sims, allowing end users to engage with a modest community of 25 agents through natural language. 
At the same time, \cite{park2022social} develops Social Simulacra, which constructs a simulated community of 1,000 personas. This system takes a designer's vision for a community—its goals, rules, and member personas—and simulates it, generating behaviors like posting, replying, and even anti-social actions.
Building on this, \cite{gao2023s} takes the concept further by constructing vast networks comprising 8,563 and 17,945 agents, respectively, designed to simulate social networks focused on the topics of Gender Discrimination and Nuclear Energy. This evolution showcases the increasing complexity and size of simulated environments in recent research.
Recent studies such as \cite{chen2023multi,kaiya2023lyfe,li2023quantifying,li2023you,ziems2023can} highlight the evolving complexity in multi-agent systems, LLM impacts on social networks, and their integration into social science research.
 
% Emotion 
% Information
\subsubsection{Gaming}
% Summary: 1. What is Game theory 2. Why this problem needs LLM-MA  3. How LLM-MA normally do in this case (Schema from section 3)
LLM-MA is well-suited for creating simulated gaming environments, allowing agents to assume various roles within games. This technology enables the development of controlled, scalable, and dynamic settings that closely mimic human interactions, making it ideal for testing a range of game theory hypotheses~\cite{mao2023alympics,xu2023exploring,gong2023mindagent}. 
%The Game theory is the study of models of strategic interaction between rational decision-makers~\cite{game_theory}. LLM-MA is suitable for game theory simulation since it can facilitate the construction of a controlled, scalable, realistic and dynamic environment of human interactions for many game theory hypothesis testing~\cite{mao2023alympics,xu2023exploring}. 
Most games simulated by LLM-MA rely heavily on natural language communication, offering a sandbox environment within different game settings for exploring or testing game theory hypotheses including reasoning, cooperation, persuasion, deception, leadership, etc. 

~\cite{akata2023playing} leverages behavioral game theory to examine LLMs' behavior in interactive social settings, particularly their performance in games like the iterated Prisoner's Dilemma and Battle of the Sexes. Furthermore, \cite{xu2023exploring} proposes a framework using ChatArena library \cite{ChatArena} for engaging LLMs in communication games like Werewolf, using retrieval and reflection on past communications for improvement, as well as the Chain-of-Thought mechanism \cite{wei2022chain}.
~\cite{light2023text} explores the potential of LLM agents in playing Resistance Avalon, introducing AVALONBENCH, a comprehensive game environment and benchmark for further developing advanced LLMs and multi-agent frameworks. ~\cite{wang2023avalon} also focuses on the capabilities of LLM Agents in dealing with misinformation in the Avalon game, proposing the Recursive Contemplation (ReCon) framework to enhance LLMs' ability to discern and counteract deceptive information. 
~\cite{xu2023language} introduces a framework combining LLMs with reinforcement learning (RL) to develop strategic language agents for the Werewolf game. It introduces a new approach to use RL policy in the case that the action and state sets are not predefined but in the natural language setting. ~\cite{mukobi2023welfare} designs the ``Welfare Diplomacy”, a general-sum variant of the zero-sum board game Diplomacy, where players must balance military conquest and domestic welfare. It also offers an open-source benchmark, aiming to help improve the cooperation ability of multi-agent AI systems. On top of that, there is a work \cite{li2023theory} in a multi-agent cooperative text game testing the agents' Theory of Mind (ToM), the ability to reason about the concealed mental states of others and is fundamental to human social interactions, collaborations, and communications. \cite{fan2023can} comprehensively assesses the capability of LLMs as rational players, and identifies the weaknesses of LLM-based Agents that even in the explicit game process, agents may still overlook or modify refined beliefs when taking actions.


\subsubsection{Psychology}
In psychological simulation studies, like in the societal simulation, multiple agents are utilized to simulate humans with various traits and thought processes. However, unlike societal simulations, one approach in psychology involves directly applying psychological experiments to these agents. This method focuses on observing and analyzing their varied behaviors through statistical methods. Here, each agent operates independently, without interacting with others, essentially representing different individuals. Another approach aligns more closely with societal simulations, where multiple agents interact and communicate with each other. In this scenario, psychological theories are applied to understand and analyze the emergent behavioral patterns. This method facilitates the study of interpersonal dynamics and group behaviors, providing insights into how individual psychological traits influence collective actions.
~\cite{ma2023understanding} explores the psychological implications and outcomes of employing LLM-based conversational agents for mental well-being support. It emphasizes the need for carefully evaluating the use of LLM-based agents in mental health applications from a psychological perspective. 
~\cite{kovavc2023socialai} introduces a tool named SocialAI school for creating interactive environments simulating social interactions. It draws from developmental psychology to understand how agents can acquire, demonstrate, and evolve social skills such as joint attention, communication, and cultural learning. 
~\cite{zhang2023exploring} explores how LLM agents, with distinct traits and thinking patterns, emulate human-like social behaviors such as conformity and majority rule. This integration of psychology into the understanding of agent collaboration offers a novel lens for examining and enhancing the mechanisms behind LLM-based multi-agents systems.
~\cite{aher2023using} introduces Turing Experiments to evaluate the extent to which large language models can simulate different aspects of human behaviors. The Turing Experiments replicate classical experiments and phenomena in psychology, economics, and sociology using a question-answering format to mimic experimental conditions. They also design a prompt that is used to simulate the responses of multiple
different individuals by varying the name. By simulating various kinds of individuals via LLM, they show that larger models replicate human behavior more faithfully, but they also reveal a hyper-accuracy distortion, especially in knowledge-based tasks. 


% Please add the following required packages to your document preamble:
% \usepackage{multirow}
% \usepackage[normalem]{ulem}
% \useunder{\uline}{\ul}{}
\begin{table*}[ht]
\centering
\resizebox{0.9\textwidth}{!}{
\begin{tabular}{l|l|l|l|l}
\hline
\textbf{Motivation} &
  \textbf{Domain} &
  \textbf{Datasets and Benchmarks} &
  \textbf{Used by} &
  \textbf{Data Link} \\ \hline
\multirow{13}{*}{Problem Solving} &
  \multirow{3}{*}{Software Development} &
  HumanEval &
  ~\cite{hong2023metagpt} & \href{https://github.com/openai/human-eval}{Link} \\
 &
   &
  MBPP &
  ~\cite{hong2023metagpt} & \href{https://github.com/google-research/google-research/tree/master/mbpp}{Link}
   \\
 &
   &
  SoftwareDev &
  ~\cite{hong2023metagpt} &
  \href{https://github.com/geekan/MetaGPT}{Link} \\ \cline{2-5} 
 &
  \multirow{4}{*}{Embodied AI} &
  RoCoBench &
  ~\cite{mandi2023roco} &
  \href{https://github.com/MandiZhao/robot-collab/tree/main/rocobench}{Link} \\
 &
   &
  Communicative Watch-And-Help (C-WAH) &
  ~\cite{zhang2023building} &
  \href{https://github.com/UMass-Foundation-Model/Co-LLM-Agents/}{Link} \\
 &
   &
  ThreeDWorld Multi-Agent Transport (TDW-MAT) &
  ~\cite{zhang2023building} &
  \href{https://github.com/UMass-Foundation-Model/Co-LLM-Agents/}{Link} \\
 &
   &
  HM3D v0.2 &
  ~\cite{yu2023conavgpt} &
  \href{https://aihabitat.org/datasets/hm3d-semantics/}{Link} \\ \cline{2-5} 
 &
  \multirow{6}{*}{Science Debate} &
  MMLU &
  ~\cite{tang2023medagents} &
  \href{https://github.com/hendrycks/test}{Link} \\
 &
   &
  MedQA &
   ~\cite{tang2023medagents} &
  \href{https://github.com/jind11/MedQA}{Link} \\
 &
   &
  PubMedQA &
  ~\cite{tang2023medagents} &
  \href{https://github.com/pubmedqa/pubmedqa}{Link} \\
 &
   &
  GSM8K &
  ~\cite{du2023improving} &
  \href{https://github.com/openai/grade-school-math}{Link} \\
 &
   &
  StrategyQA &
  ~\cite{xiong2023examining} &
  \href{https://github.com/eladsegal/strategyqa}{Link} \\
 &
   &
  Chess Move Validity &
  ~\cite{du2023improving} &
  \href{https://github.com/google/BIG-bench}{Link} \\ \hline
\multirow{15}{*}{World Simulation} &
  \multirow{3}{*}{Society} &
  SOTOPIA &
  ~\cite{zhou2023sotopia} &
  / \\
 &
   &
  Gender Discrimination &
  ~\cite{gao2023s} &
  / \\
 &
   &
  Nuclear Energy &
  ~\cite{gao2023s} &
  / \\ \cline{2-5} 
 &
  \multirow{6}{*}{Gaming} &
  Werewolf &
  ~\cite{xu2023exploring} &
  / \\
 &
   &
  Avalon &
  ~\cite{light2023text} &
  / \\
 &
   &
  Welfare Diplomacy &
  ~\cite{mukobi2023welfare} &
  / \\
 &
   &
  Layout  in the Overcooked-AI environment &
  ~\cite{agashe2023evaluating} &
  / \\
 &
   &
  Chameleon &
  ~\cite{xu2023magic} &
  \href{https://github.com/cathyxl/MAgIC}{Link} \\
 &
   &
  Undercover &
  ~\cite{xu2023magic} &
  \href{https://github.com/cathyxl/MAgIC}{Link} \\ \cline{2-5} 
 &
  \multirow{3}{*}{Psychology} &
  Ultimatum Game TE &
  ~\cite{aher2023using} &
  \href{https://github.com/microsoft/turing-experiments/tree/main}{Link} \\
 &
   &
  Garden Path TE &
  ~\cite{aher2023using} &
  \href{https://github.com/microsoft/turing-experiments/tree/main}{Link} \\
 &
   &
  Wisdom of Crowds TE &
  ~\cite{aher2023using} &
  \href{https://github.com/microsoft/turing-experiments/tree/main}{Link} \\ \cline{2-5} 
 &
  \multirow{2}{*}{Recommender System} &
  MovieLens-1M &
  ~\cite{zhang2023generative} &
  \href{https://github.com/LehengTHU/Agent4Rec/tree/master/datasets/ml-1m}{Link} \\
 &
   &
  Amazon review dataset &
  ~\cite{zhang2023agentcf} &
  / \\ \cline{2-5} 
 &
  Policy Making &
  Board Connectivity Evaluation &
  ~\cite{hua2023war} &
  \href{https://github.com/agiresearch/WarAgent}{Link} \\ \hline
\end{tabular}
}
\caption{Datasets and Benchmarks commonly used in LLM-MA studies. `` / " denotes the unavailability of  data link.}
\label{db}
\end{table*}


\subsubsection{Economy}
LLM-MA is used to simulate economic and financial trading environments mainly because it can serve as implicit computational models of humans. In these simulations, agents are provided with endowments, and information, and set with pre-defined preferences, allowing for an exploration of their actions in economic and financial contexts. This is similar to the way economists model 'homo economicus', the characterization of man in some economic theories as a rational person who pursues wealth for his own self-interest ~\cite{horton2023large}. 
There are several studies demonstrate the diverse applications of LLM-MA in simulating economic scenarios, encompassing macroeconomic activities, information marketplaces, financial trading, and virtual town simulations. Agents interact in cooperative or debate, decentralized environments. \cite{li2023large} employs LLMs for macroeconomic simulation, featuring prompt-engineering-driven agents that emulate human-like decision-making, thereby enhancing the realism of economic simulations compared to rule-based or other AI agents. \cite{anonymous2023rethinking} explores the buyer's inspection paradox in an information marketplace, reveals improved decision-making and answer quality when agents temporarily access information before purchase. \cite{li2023tradinggpt} presents an LLM-MA framework for financial trading, emphasizing a layered memory system, debate mechanisms, and individualized trading characters, thereby fortifying decision-making robustness. \cite{zhao2023competeai} utilizes LLM-based Agents to simulate a virtual town with restaurant and customer agents, yielding insights aligned with sociological and economic theories. These studies collectively illuminate the broad spectrum of applications and advancements in employing LLMs for diverse economic simulation scenarios.


\subsubsection{Recommender Systems}
% Representative works
% Motivation; How they use Multi-Agents; 
The use of the LLM-MA in recommender systems is similar to that in psychology since studies in both fields involve the consideration of extrinsic and intrinsic human factors such as cognitive processes and personality~\cite{PsychologyInformedRecommenderSystems}. One way to use LLM-MA in recommender systems is to directly introduce items to multiple LLM-based agents within diverse traits and conduct statistics of the preferences of different agents. Another way is to treat both users and items as agents and the user-item communication as interactions, simulating the preference propagation. To bridge the gap between offline metrics and real-world performance in recommendation systems, Agent4Rec~\cite{zhang2023generative} introduces a simulation platform based on LLM-MA. 1000 generative agents are initialized with the MovieLens-1M dataset to simulate complex user interactions in a recommendation environment. 
% conclusion; 
Agent4Rec shows that LLM-MA can effectively mimic real user preferences and behaviors, provide insights into phenomena like the filter bubble effect, and help uncover causal relationships in recommendation tasks. In Agent4Rec work, agents are used to simulate users and they do not communicate with each other. Different from Agent4Rec work, ~\cite{zhang2023agentcf} treats both users and items as agents, optimizing them collectively to reflect and adjust to real-world interaction disparities. This work emphasizes simulating user-item interactions and propagates preferences among agents, capturing the collaborative filtering essence.

\subsubsection{Policy Making}
Similar to simulations in gaming and economic scenarios, Policy Making requires strong decision-making capabilities to realistic and dynamic complex problems. LLM-MA can be used to simulate the policy making via simulating a virtual government or simulating the impact of various policies on different communities. These simulations provide valuable insights into how policies are formulated and their potential effects, aiding policymakers in understanding and anticipating the consequences of their decisions~\cite{policy}.
The research outlined in  \cite{xiao2023simulating} is centered on simulating a township water pollution crisis. It simulated a town located on an island including a demographic structure of different agents and township head and advisor. Within the water pollution crisis simulation, this work provides an in-depth analysis of how a virtual government entity might respond to such a public administration challenge and how information transfer in the social network in this crisis. ~\cite{hua2023war} introduces WarAgent to simulate key historical conflicts and provides insights for conflict resolution and understanding, with potential applications in preventing future international conflicts.
% The use of agent simulations in policy-making acts as a valuable roadmap for developing effective diffusion models that incorporate human reasoning and decision-making processes.


\subsubsection{Disease Propagation Simulation}
Leveraging the societal simulation capabilities of LLM-MA can also be used to simulate disease propagation.
The most recent study in \cite{williams2023epidemic} delves into the use of LLM-MA in simulating disease spread. 
% In this research, agents are uniquely programmed to autonomously reason and make decisions, facilitated through their connection to ChatGPT. 
The research showcases through various simulations how these LLM-based agents can accurately emulate human responses to disease outbreaks, including behaviors like self-quarantine and isolation during heightened case numbers. 
The collective behavior of these agents mirrors the complex patterns of multiple waves typically seen in pandemics, eventually stabilizing into an endemic state. 
Impressively, their actions contribute to the attenuation of the epidemic curve. 
% This innovative approach marks a significant stride in dynamic system modeling, introducing a sophisticated representation of human cognitive processes, reasoning, and decision-making.
~\cite{ghaffarzadegan2023generative} also discusses the epidemic propagation simulation and decomposes the simulation into two parts: the Mechanistic Model which represents the information or propagation of the virus and the Decision-Making Model which represents the agents' decision-making process when facing the virus.

%\subsubsection{More Works..}
%Because the LLM-based Multi-Agents is a fast growing research field and new applications appear almost everyday, it is difficult to include each application in the review. Hence, we maintain an updated open source repository to collect the latest work. The repository link is \href{https://github.com/taichengguo/LLM_MultiAgents_Survey_Papers}{LLM\_MultiAgents\_Survey\_Papers}.

% \subsection{Multi-Agents for Dialogue Dataset Generation}






\section{Implementation Tools and Resources}
\label{Sec4:Framework}
\subsection{Multi-Agents Framework}
% \subsection{Text-based Multi-Agents Framework}
% Agents Framework: Camel, MetaGPT, Autogen …… (xiuying



We provide a detailed introduction to three open-source multi-agent frameworks: MetaGPT~\cite{hong2023metagpt}, CAMEL~\cite{li2023camel}, and Autogen~\cite{wu2023autogen}.
They are all frameworks that utilize language models for complex task-solving with a focus on multi-agent collaboration, but they differ in their approaches and applications.

MetaGPT is designed to embed human workflow processes into the operation of language model agents, thereby reducing the hallucination problem that often arises in complex tasks. 
It does this by encoding Standard Operating Procedures into the system and using an assembly line approach to assign specific roles to different agents. 
% This structured coordination is particularly adept at deconstructing and solving complex, multi-step problems.

CAMEL, or Communicative Agent Framework, is oriented towards facilitating autonomous cooperation among agents. It uses a novel technique called inception prompting to guide conversational agents towards fulfilling tasks that are consistent with human objectives. This framework also serves as a tool for generating and studying conversational data, helping researchers understand how communicative agents behave and interact.

AutoGen is a versatile framework that allows for the creation of applications using language models. 
It is distinctive for its high level of customization, enabling developers to program agents using both natural language and code to define how these agents interact. 
This versatility enables its use in diverse fields, from technical areas such as coding and mathematics to consumer-focused sectors like entertainment.

More recently, \cite{chen2023agentverse,chen2023autoagents} introduce frameworks for dynamic multi-agent collaboration, while \cite{zhou2023agents,li2023metaagents,xie2023openagents} present platforms and libraries for building autonomous agents, emphasizing their adaptability in task-solving and social simulations.

\subsection{Datasets and Benchmarks} 
\label{Sec6:Datasets and Benchmarks}
% https://github.com/shizhl/Multi-Agent-Papers 

We summarize commonly used datasets or benchmarks for LLM-MA study in Table~\ref{db}. We observe that different research applications use different datasets and benchmarks. In the Problem solving scenarios, most datasets and benchmarks are used to evaluate the planning and reasoning capabilities by Multiple agents cooperation or debate. In World Simulation scenarios, datasets and benchmarks are used to evaluate the alignment between the simulated world and real-world or analyze the behaviors of different agents. However, in certain research applications like Science Team operations for experiments and economic modeling, there is still a need for comprehensive benchmarks.   The development of such benchmarks would greatly enhance the ability to gauge the success and applicability of LLM-MA in these complex and dynamic fields.



\input{5-Ethicals_Risks}
\section{Challenges and Opportunities}
\label{Sec7:Challenges and Opportunities}
Studies of LLM-MA frameworks and applications are 
advancing rapidly, giving rise to numerous challenges and opportunities.   We identified several critical challenges and potential areas for future study.  

%\subsection{Agents-Environment Interface: Advancing into the Multi-Modal Environment}
\subsection{Advancing into  Multi-Modal Environment}

Most previous work on LLM-MA has been focused on text-based environments, excelling in processing and generating text. However, there is a notable lack in multi-modal settings, where agents would interact with and interpret data from multiple sensory inputs and generate multiple outputs such as images, audio, video, and physical actions. Integrating LLMs into multi-modal environments presents additional challenges, such as processing diverse data types and enabling agents to understand each other and respond to more than just textual information. 

%\subsection{Agents Communication: Hallucination}
\subsection{Addressing Hallucination}
% What is hallucination
The hallucination problem is a significant challenge in LLMs and single LLM-based Agent systems. It refers to the phenomenon where the model generates text that is factually incorrect~\cite{huang2023survey}. However, this problem takes on an added layer of complexity in a multi-agent setting. In such scenarios, one agent's hallucination can have a cascading effect. This is due to the interconnected nature of multi-agent systems, where misinformation from one agent can be accepted and further propagated by others in the network. 
Therefore, detecting and mitigating hallucinations in LLM-MA is not just a crucial task but also presents a unique set of challenges. It involves not only correcting inaccuracies at the level of individual agents but also managing the flow of information between agents to prevent the spread of these inaccuracies throughout the system.

% Single Agent => Multi-Agents
% Potential challenges and solutions
%Different from the single LLM-based agent, in the Multi-Agents scenario, hallucinations can have cascading effects due to the complexity of managing multiple agents. This means that false information from one agent can propagate through the system as other agents might accept and act upon it. The hallucination may be expanded during the agent communication. Hence, detecting and mitigating the hallucination in the LLM-based Multi-Agents is a promising and challenging task.

%\subsection{Agents Capabilities Acquisition}
\subsection{Acquiring Collective Intelligence}
%As we discussed in Section~\ref{Sec3:Taxonomy}, for the feedback of the Multi-Agents Learning, % off-policy Learn from Dataset, Model
% on-policy
In traditional multi-agent systems,  agents often use reinforcement learning to learn from offline training datasets.
However, LLM-MA systems mainly learn from instant feedback, such as interactions with the environment or humans, as we discussed in Section~\ref{Sec3:Taxonomy}.
%different from Multi-Agents reinforcement learning which can learn from offline training dataset, current LLM-based Multi-Agents learning is limited to learn from instant feedback such as environment or human. 
This learning style requires a reliable interactive environment and it would be tricky to design such an interactive environment for many tasks, limiting the scalability of LLM-MA systems. Moreover, the prevailing approaches in current research involve employing Memory and Self-Evolution techniques to adjust agents based on feedback. While effective for individual agents, these methods do not fully capitalize on the potential collective intelligence of the agent network. They adjust agents in isolation, overlooking the synergistic effects that can emerge from coordinated multi-agent interactions. Hence,  jointly adjusting multiple agents and achieving optimal collective intelligence is still a critical challenge for LLM-MA.

% Effective Knowledge sharing; partial knowledge sharing
% More effective synergy 
% effective and efficient 
% trade-off between different partial knowledge sharing 


% \subsection{Scaling law of the Multi-Agents}
\subsection{Scaling Up LLM-MA Systems}
\label{scale_up}
% challenges: computational complexity; scaling law; Coordination

LLM-MA systems are composed of a number of individual LLM-based agents, posing a significant challenge of scalability regarding the number of agents. 
%single LLM-based agents, which raises a critical challenge named the scaling of the Multi-Agents.
% computational complexity
From the computational complexity perspective, each LLM-based agent, typically built on large language models like GPT-4, demands substantial computational power and memory. Scaling up the number of these agents in an LLM-MA system significantly increases  resource requirements. In scenarios with limited computational resource, it would be challenging to develop these LLM-MA systems.
%Achieving the performance of Multi-Agents in a limited computational resource would be a major challenge.

% Coordination, Communication, Learning and Adaptation
Additionally, as the number of agents in an LLM-MA system increases, additional complexities and research opportunities emerge, particularly in areas like efficient agent coordination, communication, and understanding the scaling laws of multi-agents. For instance, with more LLM-based agents, the intricacy of ensuring effective coordination and communication rises significantly. 
{As highlighted in ~\cite{multiorches2}, designing advanced Agents Orchestration methodologies is increasingly important. These methodologies aim to optimize agents workflows, task assignments tailored to different agents, and communication patterns across agents such as communication constraints between agents. Effective Agents Orchestration facilitates harmonious operation among agents, minimizing conflicts and redundancies.}
Additionally, exploring and defining the scaling laws that govern the behavior and efficiency of multi-agent systems as they grow larger remains an important area of research. These aspects highlight the need for innovative solutions to optimize LLM-MA systems, making them both effective and resource-efficient.

%\subsection{Multi-Agents Datasets and Benchmarks}
\subsection{Evaluation and Benchmarks}
% Agent Profiling Role-playing-Following ability
% Communication quality
We have summarized the datasets and benchmarks currently available for LLM-MA in  Table~\ref{db}. This is a starting point, and far from  being comprehensive. We identify two significant challenges in evaluating  LLM-MA systems and benchmarking their performance against each other. 
Firstly, as discussed in ~\cite{xu2023magic}, much of the existing research focuses on evaluating individual agents' understanding and reasoning within narrowly defined scenarios. This focus tends to overlook the broader and more complex emergent behaviors that are integral to multi-agent systems. 
Secondly, there is a notable shortfall in the development of comprehensive benchmarks across several research domains, such as Science Team for Experiment Operations, Economic analysis, and Disease propagation simulation. This gap presents an obstacle to accurately assessing and benchmarking the full capabilities of LLM-MA systems in these varied and crucial fields.
%we conclude two challenges for MultiAgents datasets and benchmarks. The first challenge is as discussed in ~\cite{xu2023magic}, most previous work focuses on evaluating the ultimate understanding and reasoning capability of agents in specific settings, overlooking the emergent capabilities in multi-agent systems. Another challenge is that in some research areas such as the Science Team for Experiment Operations, Economy, Disease propagation simulation, etc, there is still a lack of comprehensive benchmarks to evaluate the capabilities of LLM-MA.


\subsection{Applications and Beyond}
% \paragraph{Application:}
The potential of LLM-MA systems extends far beyond their current applications, holding great promise for advanced computational problem-solving in fields such as finance, education, healthcare,  environmental science, urban planning and so on. As we have discussed, LLM-MA systems possess the capability to tackle complex problems and simulate various aspects of the real world. While the current role-playing capabilities of LLMs may have limitations, ongoing advancements in LLM technology suggest a bright future. It is anticipated to have more sophisticated methodologies, applications, datasets, and benchmarks tailored for diverse research fields.
Furthermore, there are opportunities to explore LLM-MA systems from various theoretical perspectives, such as Cognitive Science~\cite{sumers2023cognitive}, Symbolic Artificial Intelligence, Cybernetics, Complex Systems, and Collective Intelligence.  Such a multi-faceted approach could  contribute to a more comprehensive understanding and innovative applications in this rapidly evolving field.

% \paragraph{Multi-Agents Systems Design and Behaviors Understanding from Diverse Perspectives:} O
% Cognitive 
% Complex Systems

\section{Conclusion}
\label{Sec8:Conclusion}
LLM-based Multi-Agents have shown inspiring collective intelligence and rapidly garnered increasing interest among researchers.
In this survey, we first systematically review the development of LLM-MA systems by positioning, differentiating, and connecting them from various aspects, regarding the agents-environment interface, the characterization of agents by LLMs, the strategies for managing agent communication and the paradigms for capability acquisition. We also summarized  LLM-MA applications for problem-solving and world simulation. By also highlighting the commonly used datasets and benchmarks and discussing challenges and future opportunities, we hope that this survey can serve as a useful resource for researchers across various research fields, inspiring future research to explore the potential of LLM-based Multi-Agents.


 

\appendix

% \section*{Acknowledgments}


%% The file named.bst is a bibliography style file for BibTeX 0.99c
\bibliographystyle{named}
\bibliography{ijcai23}

\end{document}

